% GPT-6 Astra Pro; /root/fixed_surface; 22 November 2026 UTC.
% Effort not exposed. Complete proof for independent review.
\section[HD-F1-S26: marked-divisor cohomology and the recovered Hilbert curve]{HD-F1-S26: the actual marked-divisor map is surjective in degree one, and every polarized embedded deformation recovers the covering curve}
\label{sec:f1-norm-fourfold-marked-divisors}

Work on the nonempty open of S13--S25 on which the base curve
$C$ is also nontrigonal. Thus $g(C)=5$, the connected etale
triple cover $\pi:Z\to C$ has $g(Z)=13$, and
$T$ is a specified component of the actual smooth norm fibre
in $X=\operatorname{Sym}^8 Z$. For $z\in Z$ write
\[
 X_z=\{D\in X:z\leq D\},\qquad
 D_z=X_z\cap T,\qquad M_z=\mathcal O_T(D_z).
 \tag{F1S26.1}
\]
These are the actual effective Cartier divisors and their
natural sections. For every $z$ one has
\[
 h^0(T,M_z)=1,\qquad h^1(T,M_z)=7,
 \qquad H^1(\mathcal O_T)\longrightarrow H^1(M_z)
                         \text{ is surjective}.
 \tag{F1S26.2}
\]
The kernel of the last map is the image of the tangent line
to the Abel curve at $z$ in the quotient
\[
 H^1(\mathcal O_Z)/\pi^*H^1(\mathcal O_C),
\]
under the actual restriction identification with $H^1(\mathcal O_T)$.
Consequently
\[
 h^0(D_z,\mathcal O_{D_z}(D_z))=1,
 \qquad H^1(\mathcal O_{D_z}(D_z))\longrightarrow
                                      H^2(\mathcal O_T)
                         \text{ is injective}.
 \tag{F1S26.3}
\]
The morphism $z\mapsto D_z$ identifies $Z$ with a smooth
connected component of the Hilbert scheme of $T$.

For every actual first-order deformation $(T_D,a_D:T_D\to P_D)$
arising in S25 with the given polarization on $P$ extended,
this Hilbert component extends to a smooth proper curve $Z_D$.
Its universal divisor gives an actual closed immersion
\[
                   T_D\longrightarrow\operatorname{Sym}^8 Z_D
 \tag{F1S26.4}
\]
extending the specified inclusion of $T$. The assertion here
recovers the curve and its universal divisor. The cyclic action
and the original quotient curve are addressed separately in S27.

\paragraph{1. The curve complex, its support, and the five independent parameters.}
Choose $p_0\in C$ and identify $A=\operatorname{Pic}^8C$ with
$JC$ by subtraction of $8p_0$. Put $B=\operatorname{Pic}^0(A)$.
Let $\mathcal P$ be the normalized Poincare line on $A\times B$.
The map $F:X\to A$ is the sum of the maps
$f:Z\to A$, $f(z)=\mathcal O_C(\pi z-p_0)$.
Let
\[
 \mathcal C_z=Rp_{B*}
  (\mathcal O_Z(z)\otimes(f\times1)^*\mathcal P),\qquad
 \mathcal C_0=Rp_{B*}((f\times1)^*\mathcal P).
 \tag{F1S26.5}
\]
Their perfect representatives have degrees zero and one:
the curve projection is proper and smooth, and the displayed
line bundles are flat over $B$.
Write $\varepsilon\in\operatorname{Pic}^0(C)$ for the order-three
line defining the cover and let
$\Sigma\subset B$ be the three points whose pulled-back line
on $Z$ is trivial, as specified in S21.

The degree-one line
$\mathcal O_Z(z)\otimes\pi^*\alpha$ has a section exactly at
these three parameters. Indeed a section gives
$\pi^*\alpha=\mathcal O_Z(z'-z)$ for some $z'\in Z$.
Invariance under $\sigma$ gives
\[
                         z'+\sigma z\ \sim\ z+\sigma z'.
 \tag{F1S26.6}
\]
The curve $Z$ is nonhyperelliptic on the S13 open. Thus two
distinct effective degree-two divisors cannot be linearly
equivalent: such a pencil would give a degree-at-most-two map
to $\mathbf P^1$. Equality of the two divisors in (F1S26.6)
forces $z'=z$, since $\sigma$ has no fixed point. It follows
that $\pi^*\alpha$ is trivial, giving precisely $\Sigma$.
For $\alpha\notin\Sigma$ the fibre of $\mathcal C_z$ has only
$H^1$, of dimension eleven. At a point of $\Sigma$ its fibre
dimensions are $1,12$.

Let $U=\pi^*H^1(C,\mathcal O_C)\subset H^1(Z,\mathcal O_Z)$.
It has dimension five. For the inclusion
$\mathcal O_Z\to\mathcal O_Z(z)$, the cohomology map is the
quotient by the one-dimensional boundary line
\[
 \ell_z=\operatorname{im}\left(T_zZ\longrightarrow
                                    H^1(Z,\mathcal O_Z)\right).
 \tag{F1S26.7}
\]
This boundary is the differential of the Abel map; by Serre
duality its pairing with a canonical form is evaluation at $z$.
One has $\ell_z\cap U=0$. Here are the precise evaluations
establishing this. The two nontrivial eigenspaces of canonical
forms on $Z$ are $H^0(K_C\varepsilon)$ and
$H^0(K_C\varepsilon^2)$. These systems have no base point:
Riemann--Roch says that a base point $p$ would give
$\varepsilon^j=\mathcal O_C(p-q)$ with $q\ne p$, and then
$3p\sim3q$ supplies a base-point-free degree-three pencil,
contrary to the chosen nontrigonal open. An eigenform therefore
does not vanish at $z$. Such forms annihilate the invariant
subspace $U$, so (F1S26.7) cannot lie in $U$.

At each point of $\Sigma$, cancel the invertible blocks in a
two-term locally free representative of (F1S26.5). Over the
complete regular local ring $R=\mathbf C[[t_1,\ldots,t_5]]$
the minimal representatives have ranks $1\to12$ and $1\to13$.
The derivative of their differential is the actual section
obstruction, namely cup product by $\pi^*H^1(\mathcal O_C)$.
The preceding injectivity, and the injectivity of $\pi^*$ itself,
show that these derivatives have rank five. Five entries of
each differential are consequently formal parameters, and every
other entry is in the ideal they generate. Row operations clear
the other entries. Thus the actual complexes, after completion,
are respectively
\[
 \mathcal C_z=[R\xrightarrow{(t_1,\ldots,t_5,0)}R^5\oplus R^7],
 \qquad
 \mathcal C_0=[R\xrightarrow{(t_1,\ldots,t_5,0)}R^5\oplus R^8].
 \tag{F1S26.8}
\]
The last summands modulo the maximal ideal are, intrinsically,
\[
 W_z=H^1(\mathcal O_Z(z))/U,\qquad
 V=H^1(\mathcal O_Z)/U,
 \qquad V\twoheadrightarrow W_z,\quad\ker=\overline\ell_z.
 \tag{F1S26.9}
\]
The Poincare first-order transition functions $1+\epsilon c_{ij}$
give the cup differential just used: lifting the section $1$
leaves the cocycle $c_{ij}$ in the indicated $H^1$. This fixes
the actual differential, not only its numerical rank.

\paragraph{2. The exact Fourier transform and its truncated Koszul calculation.}
Put $L_z=\mathcal O_X(X_z)$, and let
$\Phi_A=Rp_{B*}(p_A^*(-)\otimes\mathcal P)$ be the Fourier
functor. Projection formula, flat base change and the Poincare
addition isomorphism give
\[
 \Phi_A(RF_*L_z)\simeq\operatorname{Sym}^8_{\mathrm{der}}\mathcal C_z,
 \qquad
 \Phi_A(RF_*\mathcal O_X)
                  \simeq\operatorname{Sym}^8_{\mathrm{der}}\mathcal C_0.
 \tag{F1S26.10}
\]
For clarity these are the following actual maps. Pullback by the
finite quotient $q:Z^8\to X$ gives
$q^*L_z=\mathcal O_Z(z)^{\boxtimes8}$ and pulls the Poincare
line to the external product of its eight curve copies. Take
derived relative Kunneth for the product projection to $B$,
then $S_8$-invariants. The Reynolds idempotent makes invariants
exact in characteristic zero, and permutation of degree-one
factors has the Koszul sign. This is precisely the derived
symmetric power in (F1S26.10). The inclusion
$\mathcal O_X\to L_z$ is the eighth symmetric power of the
actual inclusion on the curve.

For $s=7$ or $8$, the eighth derived symmetric power of
$[R\to R^5\oplus R^s]$ in (F1S26.8) has term
$\bigwedge^j(R^5\oplus R^s)$ in degree $j$, $0\le j\le8$.
Its differential is $(8-j)t\wedge(-)$; rescaling the terms
by nonzero constants gives the usual wedge Koszul differential.
In degrees through seven, including the outgoing degree-seven
differential, this is the full cochain Koszul complex for the
five parameters tensored with $\bigwedge^\bullet R^s$ with
zero differential. The full five-parameter complex resolves
$k[-5]$, so its truncation through degree seven is
\[
 k[-5]\ \oplus\ k^s[-6]\ \oplus\ k^{\binom s2}[-7].
 \tag{F1S26.11}
\]
This is an isomorphism of the actual truncated complexes in
the derived category, not just a list of their cohomology
dimensions. Outside $\Sigma$ the curve complex is a vector
bundle placed in degree one, so its eighth symmetric power
has cohomology only in degree eight. Thus globally the
truncation in (F1S26.11) is the sum of these skyscraper
complexes over the three points of $\Sigma$. To justify the
passage from completion, all these cohomology sheaves have
finite support. Their extension groups are the direct sums
of the local extension groups at those points, and completion
is an isomorphism on these finite-length extension groups.
The zero local Postnikov classes in (F1S26.11) are therefore
zero globally.

Let $G_s$ denote the remaining coherent cohomology sheaf in
degree eight. The exact truncation triangle is consequently
\[
 \bigoplus_{\alpha\in\Sigma}
       (k_\alpha[-5]\oplus k_\alpha^s[-6]
                         \oplus k_\alpha^{\binom s2}[-7])
 \longrightarrow\Phi_A(\mathcal B_s)
 \longrightarrow G_s[-8]\longrightarrow,
 \tag{F1S26.12}
\]
where $\mathcal B_7=RF_*L_z$ and
$\mathcal B_8=RF_*\mathcal O_X$.

Mukai inversion in dimension five gives
$\Psi=( -1_A)^*\Phi_B[5]$ as inverse to $\Phi_A$.
In particular $\Psi(k_\alpha)=\lambda_\alpha[5]$ for the
specified flat line $\lambda_\alpha$ on $A$; changing the
Poincare sign replaces $\Sigma$ by its inverse, the same set.
Applying $\Psi$ to (F1S26.12) gives a triangle whose first
term is
\[
 \bigoplus_{\alpha\in\Sigma}
  (\lambda_\alpha\oplus\lambda_\alpha^s[-1]
                         \oplus\lambda_\alpha^{\binom s2}[-2]),
 \tag{F1S26.13}
\]
and whose last term is $( -1_A)^*\Phi_B(G_s)[-3]$.
The latter has derived fibres in degrees at least three.
Indeed the sheaf $p_B^*G_s\otimes\mathcal P$ on $B\times A$
is flat over $A$, and proper flat base change identifies its
derived fibre after pushforward with
$R\Gamma(B,G_s\otimes\mathcal P_a)$, in degrees $0$ through
$5$. The shift $[-3]$ puts these in degrees $3$ through $8$.
This proves the required Tor-amplitude assertion even when
$G_s$ is not locally free.

It follows in particular that the derived fibre of
$RF_*L_z$ at the specified canonical point has dimensions
$3,21$ in degrees zero and one. For $RF_*\mathcal O_X$ they
are $3,24$. The map between their degree-one fibres is,
at each Fourier support point, the quotient $V\to W_z$ in
(F1S26.9). To see this without choosing a compatibility of
all splittings, apply the actual map of curve complexes to
the Koszul computation. On its degree-six cohomology it is
$\det(U)\otimes V\to\det(U)\otimes W_z$, with the common
determinant factor and the quotient of (F1S26.9). Off-diagonal
extension components between different cohomological degrees
of (F1S26.13) become positive Ext classes between flat lines.
Their pullbacks to a point are zero. Therefore the asserted
degree-one fibre map is exactly the displayed quotient and
is surjective of total rank twenty-one.

\paragraph{3. Actual excess correction and the three component idempotents.}
The derived fibre of $F$ has the exact S21 triangle with
excess line $\mathcal O_{\mathcal T}(-h)$, where
$\mathcal T=T_0\amalg T_1\amalg T_2$ and
$h=p^*\mathcal O_{|K_C|}(1)$. Tensor it by $L_z$. The
actual norm of a marked-point divisor gives the line identity
\[
 \mathcal O_{\mathcal T}(h)
  =M_z\otimes M_{\sigma z}\otimes M_{\sigma^2z},
 \qquad
 M_z\otimes\mathcal O_{\mathcal T}(-h)
       =(M_{\sigma z}\otimes M_{\sigma^2z})^{-1}.
 \tag{F1S26.14}
\]
The first equality can be checked after $Z^8\to X$: the
equation for a canonical divisor to contain $\pi z$ pulls
back to the product of the equations for its eight points
to equal one of $z,\sigma z,\sigma^2z$, including repeated
points with their scheme multiplicities.
Each point-divisor line on $X$ is ample: its pullback to
$Z^8$ is the external product of eight ample degree-one
lines, and ampleness descends along a finite surjection.
Thus the last line of (F1S26.14) is antiample on each smooth
fourfold. Kodaira vanishing and Serre duality give its
$H^q=0$ for $q<4$. The exact excess triangle therefore
identifies the derived and ordinary cohomology in degrees
zero and one. In degree two it gives only an injection,
with a possible boundary to $H^4$ of that antiample line;
no isomorphism in degree two is asserted here.

We now pass to individual components, retaining the actual
module structure. Choose a trivialization of $\pi^*\varepsilon$
whose cube is one relative to the chosen trivialization of
$\varepsilon^3$. Its external eighth power gives the three
compatible trivializations of $F^*\alpha$, $\alpha\in\Sigma$.
They act both on $RF_*\mathcal O_X$ and on $RF_*L_z$.
On the canonical fibre their algebra is
$H^0(\mathcal T,\mathcal O_{\mathcal T})=\mathbf C^3$;
its three primitive idempotents are precisely the three
actual component projections. Under (F1S26.10) the action
is translation of the Fourier support points $\Sigma$,
with their coherent multiplication, so it permutes the
three equal blocks in (F1S26.11) transitively. A representation
with these three blocks and this order-three action is the
regular representation tensored with a block: choose one
block and transport its basis around the order-three orbit.
Its three character spaces have the same dimension.
Consequently (F1S26.13) and (F1S26.14) give, on each component,
\[
 h^0(M_z)=1,\quad h^1(M_z)=7,
 \quad H^1(\mathcal O_T)=V\twoheadrightarrow W_z=H^1(M_z).
 \tag{F1S26.15}
\]
The latter quotient is compatible with the component
idempotents because it is induced by the actual section
inclusion. The usual restriction map from $H^1(\mathcal O_Z)$
to the symmetric product and then to $T$ gives the stated
identification with $V$: its invariant subspace is killed
by the constant norm map, and its quotient is the S21
Albanese pullback isomorphism. This proves (F1S26.2),
including the kernel line.

\paragraph{4. The actual Hilbert curve.}
The canonical section of $M_z$ is nonzero on every component:
choose a canonical divisor on $C$ avoiding $\pi z$ and a
point of the corresponding sheet component above it.
Thus $H^0(\mathcal O_T)\to H^0(M_z)$ is an isomorphism.
The actual Cartier-divisor sequence
\[
 0\longrightarrow\mathcal O_T\longrightarrow M_z
       \longrightarrow\mathcal O_{D_z}(D_z)\longrightarrow0
 \tag{F1S26.16}
\]
and (F1S26.15) prove (F1S26.3). The restriction of the
universal divisor in $Z\times X$ to $Z\times T$ is a
flat family of effective Cartier divisors on $T$: its
section is a nonzerodivisor on each fibre, so the quotient
is flat over the smooth curve parameter. It defines
$Z\to\operatorname{Hilb}(T)$. Its differential is nonzero
at every point, because its image under the boundary of
(F1S26.16) is the nonzero projected line
$\overline\ell_z$ in (F1S26.9).

This morphism is injective. The line difference of two
point divisors on $X$ is the Abel pullback of the line in
$\operatorname{Pic}^0(JZ)$ corresponding to $\mathcal O_Z(z-z')$.
On $T$, S21 identifies $\operatorname{Pic}^0(T)$ with
$\operatorname{Pic}^0(P)$. The dual of the actual inclusion
$P\hookrightarrow JZ$ has kernel the connected subvariety
$\pi^*JC$. Hence $M_z=M_{z'}$ would imply
$\mathcal O_Z(z-z')\in\pi^*JC$, which forces $z=z'$ by
the degree-two argument (F1S26.6). Equality of divisors
implies equality of their lines, proving injectivity.
The proper unramified injective map from the smooth curve
$Z$ is a closed immersion into its image. Its target
Hilbert tangent space has dimension one by (F1S26.3), and
the image has dimension one. Thus the Hilbert scheme is
regular of dimension one at every point of the image;
the image is open as well as closed in its component.
It is exactly that smooth connected component.

\paragraph{5. The exact line identity needed for polarized deformation.}
We record why the relevant line classes extend, without
replacing a line identity by a rational numerical equality.
On $X$ the known canonical class formula is
$c_1(\omega_X)=\theta+4\eta$ (S13). For any chosen theta
line $\Theta$ on $JZ$, put
$R_z=\omega_X\otimes u^*\Theta^{-1}\otimes L_z^{-4}$,
where $u:X\to JZ$ is the chosen Abel map.
Its first Chern class is zero rationally. Pull back by
$q:Z^8\to X$. The ramification formula is the actual
line equality
\[
 q^*\omega_X=K_Z^{\boxtimes8}
                      \otimes\mathcal O_{Z^8}(-\sum_{i<j}\Delta_{ij}).
 \tag{F1S26.17}
\]
The integral second cohomology of $Z^8$ is torsion free
by the integral Kunneth formula for curves. Hence
$c_1(q^*R_z)=0$ integrally, and $q^*R_z\in\operatorname{Pic}^0(Z^8)$.
The latter Picard variety is the product of the eight
curve Picard varieties. Permutation invariance forces
$q^*R_z$ to be the external product of eight copies of
one degree-zero line. The Abel pullback isomorphism
$\operatorname{Pic}^0(JZ)\to\operatorname{Pic}^0(Z)$ and
the Poincare addition map therefore produce a specified
$\gamma_z'\in\operatorname{Pic}^0(JZ)$ with
$q^*R_z=q^*u^*\gamma_z'$.

Pullback by $q$ is injective on line bundles. Indeed a
trivialization of a trivial pullback has a permutation
character, since all global units on the proper connected
$Z^8$ are constant. On the small diagonal every permutation
fixes the point and acts trivially on the fibre of a
pulled-back line; the character is consequently trivial.
The trivialization descends by invariants. It follows that
$R_z=u^*\gamma_z'$ as actual lines.
The same argument, or the normalized Poincare addition
isomorphism directly, shows that each ratio
$L_{\sigma^iz}L_z^{-1}$ is pulled back from
$\operatorname{Pic}^0(JZ)$.

Adjunction, $\det N_{T/X}=\mathcal O_T(-h)$ from S13,
and (F1S26.14) now give the actual identity
\[
 \omega_T\otimes a^*\mathcal B^{-1}
                         =M_z\otimes a^*\gamma_z,
            \qquad \gamma_z\in\operatorname{Pic}^0(P),
 \tag{F1S26.18}
\]
where $\mathcal B$ is any fixed representative of the
given polarization $\beta$ on $P$. The change from the
chosen translated restriction of $\Theta$ to $\mathcal B$
is another degree-zero line, included in $\gamma_z$.
Every factor and correction in (F1S26.18) has thus been
specified by actual pullback maps.

For a polarized deformation $(T_D,a_D,P_D)$, the relative
canonical bundle and the representative of $\beta$ extend.
Every $\gamma_z\in\operatorname{Pic}^0(P)$ lifts to the
smooth relative dual abelian scheme over $D$.
Equation (F1S26.18) therefore extends $M_z$ to a line on
$T_D$. Its section obstruction is an element of $H^1(M_z)$.
Changing this line extension by an infinitesimal degree-zero
line with class $v\in H^1(\mathcal O_T)$ changes the
section obstruction by the image of $v$ under the actual
map in (F1S26.15): on Cech overlaps this is multiplication
of the old section by the transition $1+\epsilon v_{ij}$.
Surjectivity cancels the obstruction. Thus the section
lifts, and its zero divisor is flat over $D$ since its
central section is a nonzerodivisor. Every point $D_z$ of
the intrinsic Hilbert curve has an actual relative lift.

\paragraph{6. The whole Hilbert curve and its universal divisor lift.}
The relative Hilbert scheme exists here: $\omega_T$ is
ample by S13, hence $\omega_{T_D/D}$ is relatively ample
over the nilpotent base and $T_D$ is projective over $D$.
Take the Hilbert polynomial of the divisors $D_z$.
Its central Hilbert scheme is smooth of dimension one
along the proper open-and-closed component $Z$ just
proved. The corresponding open-and-closed part of the
relative Hilbert scheme is smooth over $D$.

Here is a local algebraic verification of the last claim,
including flatness. Embed a neighbourhood of that part
in affine space over $D$. Lift a set of central equations
with independent Jacobian rows cutting out the smooth
central curve; their zero set is a smooth relative curve
$U_D$. The remaining ideal of the Hilbert scheme inside
$U_D$ is contained in $\epsilon\mathcal O_{U_D}$, and
can be written $\epsilon\mathcal J$ for an ideal
$\mathcal J\subset\mathcal O_U$ on the reduced central
curve. A central point admits a $D$-valued lift precisely
when every local member of $\mathcal J$ vanishes at that
point: evaluation of $\epsilon f$ is $\epsilon f(x)$,
independent of the infinitesimal displacement of the point.
Paragraph5 supplies such a lift for every point of $U$.
The ideal $\mathcal J$ therefore vanishes on the reduced
finite-type curve and is zero. The Hilbert scheme equals
$U_D$ there. This proves both relative smoothness and
flatness, and constructs the smooth proper curve $Z_D$.

Let $\mathcal D_D\subset Z_D\times_D T_D$ be its actual
universal divisor. Its projection to $T_D$ reduces to
the degree-eight universal divisor on $Z\times T$.
It is proper and has finite fibres, hence is finite.
The universal family is flat over $Z_D$, hence over $D$;
its central fibre is flat of degree eight over $T$.
The fibrewise flatness criterion, with $T_D$ flat over
$D$, gives flatness over $T_D$ and rank eight. Thus it
is an actual family of effective degree-eight divisors
on $Z_D$, and represents the morphism (F1S26.4) to the
relative symmetric product. The morphism is proper and
reduces to the original closed immersion. It is finite
(proper and quasi-finite), and its map of finite algebras
is surjective by Nakayama after reduction. It is therefore
a closed immersion. This proves all asserted deformation
maps without assuming that deformations of $T$ come from
deformations of the original cover.

\paragraph{7. Compatibility of the additional open with the period certificate.}
Take the same connected component of the fine cyclic-cover
space on which S23 supplies its nonempty period-rank-twelve
open. The nonzero order-three line cover is finite etale
over fine genus-five curve moduli. Each connected component
is smooth and hence irreducible, and its finite etale image
is both open and closed in the corresponding connected
base component. Thus it maps onto that component. The
trigonal locus in genus five is closed: it is the proper
image of the relative $W^1_3$ in the degree-three Picard
scheme. Its dimension is at most eleven. Indeed a degree-three
map has total ramification degree fourteen, so its branch
values vary in dimension at most fourteen; quotient by
$\operatorname{PGL}_2$ leaves at most eleven, and fixed
branch data have discrete Hurwitz fibres. The curve-moduli
component has dimension twelve. The nontrigonal inverse
image is consequently dense open in the same irreducible
cover component and meets the existing S13--S23 open.
This argument does not require the entire cyclic-cover
space to be irreducible.

\paragraph{Primary source and proof scope.}
Mukai, \emph{Duality between $D(X)$ and $D(\widehat X)$
with its application to Picard sheaves}, Nagoya Math.
J.~81 (1981),153--175, DOI10.1017/S002776300001922X,
Proposition1.3, Theorem2.2 and its full proof, Example2.6,
and formulas(3.1),(3.7) with proofs were read in the
primary publisher PDF. The inversion proof computes the
actual composition kernel by Poincare addition, obtaining
the graph of $-1$ shifted by $-5$; these are the shifts
used above. The local curve model and its symmetric-power
Koszul calculation are proved here, as is the Tor-amplitude
of the remaining inverse-transform term. The excess
triangle and all component maps are the exact S21 maps;
only degrees zero and one are identified with ordinary
cohomology. The degree-two boundary is retained explicitly.
All statements concern the actual genus-five triple-cover
family on the indicated open and its actual divisors.
